\chapter*{结语}
\addcontentsline{toc}{chapter}{结语}

本书有意选择了一条有所取舍的路径，来梳理生成建模的数学基础。我们从线性表示开始：把矩阵理解为变换，把 PCA 理解为保持方差的投影，并把自编码器理解为基于重构的压缩。随后，概率主成分分析把这一确定性图景转化为潜变量生成模型，并由此引入先验、似然、边缘似然和后验这些核心建模对象。

变分自编码器把这一潜变量视角推广到了非线性情形，此时精确的边缘似然和后验已不再可得。这自然引出了证据下界、变分推断、近似后验以及重参数化技巧。扩散模型则进一步把生成过程从单一潜变量转变为一整条带噪潜在状态轨迹，使生成成为逐步加噪过程的一个学习得到的逆过程。

连续时间和基于得分的章节从第二种视角考察了同一个广义问题。我们不再只把生成描述为离散链，而是研究运动、密度演化、布朗噪声、Liouville 方程、Fokker--Planck 方程、Langevin 动力学、得分匹配以及逆时间扩散。这说明，生成建模也可以被理解为学习概率空间中的运动方向。

后续章节考察了两条互补路径。归一化流和自回归模型展示了如何通过结构设计保留精确似然：要么使用具有易于计算的 Jacobian 的可逆变换，要么采用有序的条件分解。GAN、Wasserstein GAN 和能量模型则说明，生成学习也可以超越易于计算的似然，转而依赖比较、几何差异或标量能量景观。

综合来看，这些章节把生成式人工智能呈现为一组反复出现的数学思想，而不是一堆彼此无关的架构。这些思想包括：表示、变换、边缘化、近似、加噪与去噪、密度演化、得分场、似然构造、分布比较以及能量塑形。它们也仍在不断涌现的新型生成建模框架中反复出现。

例如，流匹配（flow matching）~\cite{lipman2023flow} 和整流流（rectified flows）~\cite{liu2022rectified} 可以看作连续时间输运视角的进一步发展：它们学习一个向量场，将概率质量从简单分布输运到数据分布。一致性模型（consistency models）~\cite{song2023consistency} 则可看作对扩散模型采样效率问题的另一种回应，其目标是在减少采样步数的同时，尽可能保留迭代去噪所带来的生成质量。

因此，本书的目的，是建立一套紧凑的数学基础，使读者能够更加从容地理解这些后续发展。一旦掌握潜变量、变分界、扩散路径、密度演化、得分场、精确似然模型、对抗式比较以及能量景观，许多现代扩展方法就更容易被放置在同一张概念导图中理解。
